\chapter{Conclusiones}

\noindent Este capítulo sintetiza las aportaciones del trabajo, revisa el grado de cumplimiento de los objetivos planteados y delimita las tareas necesarias para completar la evaluación experimental independiente.

\section{Conclusiones principales}

El trabajo confirma que la identificación de una posible compraventa en el mercado de artículos de \textit{Counter-Strike 2} no puede reducirse a predecir una subida o una bajada de precio. La decisión depende de los precios de entrada y salida, las comisiones, la conversión de moneda, el periodo de bloqueo de intercambio, el capital disponible y la exposición acumulada de la cartera. Por ello, el resultado principal no es una regla aislada de compra, sino un modelo que permite estudiar decisiones secuenciales con esas condiciones de forma reproducible.

La propuesta se ha formalizado como un modelo MARL aplicado al caso de estudio. El estado global combina información de mercado, estado de cartera y señal supervisada; cada agente recibe una observación local coherente con su responsabilidad. \textit{Scout} identifica candidatos, \textit{Trader} propone mantener, comprar o vender, y \textit{Portfolio} aprueba o rechaza la propuesta según la información de cartera. El simulador interpreta la acción conjunta, comprueba los límites definidos y actualiza la cartera simulada. Esta separación permite entrenar con información global para evaluar la coordinación y ejecutar cada política con su observación local.

Sobre ese modelo se ha construido un flujo completo de adquisición, normalización y persistencia de datos, elaboración de conjuntos temporales, estimación supervisada, simulación económica y consulta mediante una consola local. El conjunto \texttt{trading\_profit\_v2} contiene operaciones rentables y no rentables en entrenamiento, validación y prueba. Además, el entorno MARL ha ejecutado controles de compra, bloqueo y venta, y una matriz de 21 entrenamientos con 3.552 episodios sobre el bloque de validación. Estas evidencias demuestran que los componentes se integran y pueden evaluarse con una configuración trazable.

No obstante, los resultados no permiten concluir que la arquitectura multiagente obtenga una rentabilidad superior a una estrategia sencilla. Las diferencias observadas entre configuraciones son pequeñas y todavía no se ha realizado la evaluación independiente sobre el conjunto de prueba. La contribución del trabajo consiste, por tanto, en proporcionar una base técnica y metodológica para realizar esa comparación sin confundir una señal, una simulación o una comprobación funcional con una operación rentable garantizada.

\section{Cumplimiento de objetivos}

El objetivo general de diseñar y evaluar mediante simulación una arquitectura multiagente para apoyar la identificación de operaciones en mercados de activos digitales se ha cumplido de forma parcial: el diseño, la implementación y la evaluación inicial están disponibles, mientras que la comparación económica independiente queda pendiente. El estado de los objetivos específicos es el siguiente:

\begin{itemize}
  \item \textbf{Análisis del caso de estudio: cumplido.} Se ha caracterizado el mercado de artículos de \textit{Counter-Strike 2}, sus plataformas, los costes de transacción, las divisas, el periodo de bloqueo de intercambio y las particularidades que diferencian este mercado de uno bursátil.
  \item \textbf{Construcción de datos: cumplido con limitaciones de cobertura.} El sistema adquiere, normaliza y conserva datos de las fuentes disponibles con su fecha, moneda y origen. Los cortes construidos documentan tanto las primeras limitaciones como el conjunto temporal actual utilizado para el entrenamiento.
  \item \textbf{Simulación y estimación: cumplido de forma inicial.} Se han convertido las reglas económicas en funciones comprobables, se han construido conjuntos de compraventa y se han entrenado exploraciones supervisadas. La señal resultante se utiliza como información auxiliar de decisión, no como una orden de compra.
  \item \textbf{Coordinación multiagente: cumplido en formulación, entorno y entrenamiento inicial.} Se han definido el estado, las observaciones, las acciones, las recompensas y las restricciones de los tres roles. El entorno ejecuta episodios y la matriz MARL se ha entrenado con semillas y validación, pero falta la comparación final con referencias homogéneas.
\end{itemize}

\section{Limitaciones}

La principal limitación experimental sigue siendo la distribución de los datos. Aunque \texttt{trading\_profit\_v2} contiene las dos clases en las tres particiones temporales, las operaciones rentables predominan en validación y prueba. Además, muchos artículos aparecen tanto en entrenamiento como en prueba. La separación por fecha evita utilizar información futura, pero no demuestra que las estimaciones generalicen a artículos que no se hayan visto previamente.

También existe una limitación operativa. Las operaciones se evalúan con precios históricos y reglas económicas configuradas; no se envían órdenes reales, no se transfieren artículos y no se utiliza capital real. La sincronización entre la entrada en \acroref{BUFF}{BUFF} y la salida estimada en Steam requiere más capturas alineadas para aproximar mejor las condiciones reales de ejecución. Asimismo, el \acroref{OCR}{OCR} está validado a nivel funcional, pero todavía no cuenta con un conjunto de imágenes etiquetado para medir su exactitud.

Por último, la matriz MARL confirma que el entrenamiento, la parada temprana y la validación se ejecutan de forma reproducible, pero no demuestra todavía una ventaja robusta de una configuración sobre otra. La prueba aislada no se ha utilizado para seleccionar parámetros y debe conservarse para la comparación final. Estas limitaciones acotan las conclusiones a viabilidad técnica, coherencia metodológica y evidencia experimental inicial.

\section{Trabajo futuro}

La prioridad inmediata es evaluar una sola vez en la prueba aislada la configuración seleccionada mediante validación. Esta evaluación debe compararla con referencias homogéneas: mantener el efectivo, comprar cuando el margen neto observado sea positivo y un agente único sujeto a las mismas restricciones. Debe comunicar el retorno de cartera, el beneficio neto, las operaciones cerradas, las posiciones abiertas y los rechazos por restricciones como media e intervalo entre semillas.

En paralelo, será necesario ampliar la cobertura temporal y el número de artículos de los históricos de mercado, así como conservar más observaciones sincronizadas entre plataformas. Esto permitirá reducir el desequilibrio de etiquetas, evaluar la generalización a artículos no vistos y comprobar si las variables históricas o la señal supervisada aportan valor de forma estable. Para el OCR, el siguiente paso es construir un conjunto de capturas etiquetadas que permita medir por separado el reconocimiento de artículo, calidad, plataforma, precio y moneda.

Finalmente, una vez cerrada la comparación básica, podrán estudiarse variantes de la arquitectura: modificar la combinación entre recompensa cooperativa e individual, retirar la señal supervisada o comparar los tres roles con una política única. Estas extensiones solo deben considerarse mejoras si se seleccionan con validación y superan las mismas referencias en la prueba independiente.
